\documentclass{article}
\usepackage{amsmath,amssymb,xurl}
\usepackage[hidelinks]{hyperref}
% Shared commands for notes in docs/paper-gaps/.

\DeclareUnicodeCharacter{2080}{\ifmmode{}_{0}\else\textsubscript{0}\fi}
\DeclareUnicodeCharacter{2081}{\ifmmode{}_{1}\else\textsubscript{1}\fi}
\DeclareUnicodeCharacter{2082}{\ifmmode{}_{2}\else\textsubscript{2}\fi}
\DeclareUnicodeCharacter{2099}{\ifmmode{}_{n}\else\textsubscript{n}\fi}

\newcommand{\Lean}{\textsc{Lean}}
\ExplSyntaxOn
\NewDocumentCommand{\leanid}{m}{
  \ifmmode
    \textsf{\detokenize{#1}}
  \else
    \group_begin:
    \urlstyle{sf}
    \tl_set:Nn \l_tmpa_tl {#1}
    \tl_replace_all:Nnn \l_tmpa_tl {\_} {_}
    \exp_args:NV \path \l_tmpa_tl
    \group_end:
  \fi
}
\ExplSyntaxOff
\providecommand{\tr}{\operatorname{tr}}
\newcommand{\tnleanIssueBase}{https://github.com/LionSR/TNLean/issues/}
\newcommand{\tnleanPrBase}{https://github.com/LionSR/TNLean/pull/}
\newcommand{\tnleanDocBase}{https://sirui-lu.com/TNLean/docs/}
\newcommand{\ghissue}[1]{\href{\tnleanIssueBase#1}{GitHub issue \##1}}
\newcommand{\ghpr}[1]{\href{\tnleanPrBase#1}{PR \##1}}
\newcommand{\doclink}[1]{\href{\tnleanDocBase#1.html}{\path{#1}}}

% Dirac notation and the identity operator, as in blueprint/src/macros/common.tex.
% \providecommand keeps note-local definitions authoritative.
\usepackage{dsfont}
\providecommand{\ket}[1]{|#1\rangle}
\providecommand{\bra}[1]{\langle #1|}
\providecommand{\braket}[2]{\langle #1 | #2 \rangle}
\providecommand{\ketbra}[2]{|#1\rangle\!\langle #2|}
\providecommand{\Id}{\mathds{1}}
\providecommand{\one}{\mathds{1}}
\providecommand{\C}{\mathbb{C}}
\providecommand{\MN}[1]{M_{#1}(\C)}
\providecommand{\MD}{M_D(\C)}

% Verdict marker: \gapnote{<kind>}{<status>}. Kind is the policy
% classification (clarification, local-correction, scope-restriction,
% unfaithful, false-source, open-gap); status is open, wip (elimination
% underway), resolved, or historical. Machine-read by the site index;
% prints a verdict line.
\newcommand{\gapnote}[2]{\par\noindent\textbf{Verdict:} #1 (#2).\par}

\title{Flux Vacuum Density inside an Admissible Physical Envelope}
\author{}
\date{}
\begin{document}
\maketitle
\gapnote{scope-restriction}{open}

Schuch, Cirac, and P\'erez-Garc\'ia,
\href{https://arxiv.org/abs/1001.3807v3}{arXiv:1001.3807v3}~\cite{Schuch2010PEPS},
Lemma 6.14, source
\path{Papers/1001.3807/paper_v3.tex}, lines 2199--2214,
compare the local physics away from flux endpoints with the vacuum.
An endpoint-free annulus can detect an enclosed endpoint, so absence of
endpoints from the observed sites alone is insufficient.

The theorem
\leanid{TNLean.PEPS.IsGIsometric.torusFlux_normalized_density_eq_vacuum_of_subset_isSimplyConnected}
extends the explicit native density theorem to every observation $R$
contained in an admissible physical envelope $S$. Only $S$ is required
to have connected induced graph and simply connected closed-cell realization.
Its computed four-corner forbidden dual set must separately have simply
connected realization, and both path endpoints must lie outside that set.
Both periods are at least three and the tensor is regular group-isometric
for a finite group. Noncentral insertion elements and arbitrary winding
are allowed. Every hypothesis of the native envelope theorem is preserved.

The smaller region $R$ can be disconnected, have holes or be empty.
Its normalized density is obtained by the actual partial trace of the
derived density on $S$. Trace preservation gives the same normalization
scalar, and nonvanishing of the coefficient vector proves both cuts on $R$
nonzero. No marginal equality or nonzero-cut witness is supplied.
The generic matrix statement is
\leanid{TNLean.PEPS.regionPhysicalCutMatrix_normalized_reduced_eq_of_subset}.

This sufficient condition does not classify all admissible observations.
For an annulus surrounding an endpoint, filling its hole into an envelope
can include that endpoint in the forbidden dual set. The envelope hypotheses
then fail, consistently with the Wilson-loop obstruction. No theorem here
asserts that endpoint exclusion alone gives an admissible envelope or
classifies restricted homotopy. Unnormalized reduced matrices and global
winding states are not equated with vacuum. The broader source wording
remains open under \ghissue{8703}; \ghissue{8947} is unchanged.
\bibliographystyle{alpha}
\bibliography{references}
\end{document}
